\documentclass[12pt]{article}
\usepackage[utf8]{inputenc}
\usepackage[T1]{fontenc}
\usepackage{hyperref}
\title{Sambe Language Resources for K-12 Curriculum Development}
\author{Compiled by Claude Sonnet 5.5 (claude-sonnet-5-5) for Andr\'as Kornai}
\date{}
\begin{document}
\maketitle

\section{Introduction}
Sambe (ISO 639-3 xab; Glottocode samb1307) is a Plateau language of Nigeria. It was spoken in a single village in Kaduna State. By 2005 only two very elderly speakers remained, and the language is presumed to be no longer spoken. The community keeps its Sambe identity but speaks a dialect of Ninzo. Sambe is linguistically notable for contrasting the labial-velar stops /k\textsuperscript{\ }p/ and a labialized counterpart, a rare distinction. Its closest relative is Hasha, within the little-known Alumic group. The search found no school materials, primers or native-language texts. A curriculum would have to be built from the scholarly documentation below.

\section{Primary Educational Resources}
No primers, textbooks or school programmes for Sambe turned up in the search. Any classroom use would need to be adapted from the wordlist in the documentation listed in the next section.

\section{Linguistic and Reference Materials}
\begin{itemize}
\item Roger Blench, ``Final Records of the Sambe Language of Central Nigeria: Phonology, noun morphology, and wordlist'', \textit{Language Documentation \& Conservation} 9 (2015), pp.\ 192--228. It is based on fieldwork in 2001 and 2005 and is the main body of data. \url{https://scholarspace.manoa.hawaii.edu/items/2065949b-c8ec-431b-bc2c-07995aa11e54}
\item Roger Blench and Barau Kato (2006), ``The Sambe language of Central Nigeria and its affinities'', listed by the Endangered Languages Project. \url{https://endangeredlanguages.com/node/95664}
\item Endangered Languages Project entry for Sambe, which rates it critically endangered. \url{https://endangeredlanguages.com/lang/8949}
\item Glottolog reference record for Blench 2015. \url{https://glottolog.org/resource/reference/id/474792}
\item Wikipedia overview of the Sambe language. \url{https://en.wikipedia.org/wiki/Sambe_language}
\end{itemize}

\section{Religious and Community Resources}
No Sambe-language religious texts or community organisations were found in the search results.

\section{Digital and Audio Resources}
No audio or video recordings, corpora or software were found. The Foundation for Endangered Languages bulletin carries a notice about the language: \url{https://www.ogmios.org/ogmios_files/206.htm}

\section{Curriculum Development Context}
Sambe has effectively no living speakers, so any curriculum would be a heritage and cultural-awareness programme and not an immersion one. The documented wordlist and phonology could support lessons on local history, on the community's shift to Ninzo, and on language endangerment. Blench's published data is the only substantial source found. Further sources, and any contact with the Sambe community, would need to be established through the authors of the documentation above.

\end{document}